%% ============================================================
%%  SECCIÓN 7 — Resultados Experimentales
%%  Responsable: Minilux
%% ============================================================

\section{Resultados Experimentales}
\label{sec:resultados}

En esta sección se presentan los datos obtenidos durante las mediciones
en el laboratorio. Se reportan primero los parámetros geométricos del
anillo y los jinetillos empleados en el primer método, y a continuación
los datos del segundo método.

%% ------------------------------------------------------------------
\subsection{Primer método --- Balanza de Mohr--Westphal}
%% ------------------------------------------------------------------

Los parámetros geométricos del anillo metálico se presentan en la
Tabla~\ref{tab:param_anillo}.

\begin{table}[H]
    \centering
    \caption{Parámetros geométricos del anillo metálico.}
    \label{tab:param_anillo}
    \begin{tabular}{lcc}
        \toprule
        \textbf{Magnitud} & \textbf{Símbolo} & \textbf{Valor} \\
        \midrule
        Radio externo        & $R_{\text{ext}}$ & $(5{,}040 \pm 0{,}005)\ \text{cm}$ \\
        Radio interno        & $R_{\text{int}}$ & $(4{,}980 \pm 0{,}005)\ \text{cm}$ \\
        Longitud de contacto & $L$              & \qty{62{,}957e-2}{\meter}           \\
        \bottomrule
    \end{tabular}
\end{table}

Las masas de los jinetillos disponibles, medidas con la balanza de
precisión, se consignan en la Tabla~\ref{tab:jinetillos}.

\begin{table}[H]
    \centering
    \caption{Masas de los jinetillos empleados.}
    \label{tab:jinetillos}
    \begin{tabular}{cc}
        \toprule
        \textbf{Jinetillo} & \textbf{Masa (g)} \\
        \midrule
        $J_1$ & $20{,}2 \pm 0{,}1$ \\
        $J_2$ & $18{,}7 \pm 0{,}1$ \\
        $J_3$ & $22{,}2 \pm 0{,}1$ \\
        $J_4$ & $ 1{,}2 \pm 0{,}1$ \\
        $J_5$ & $17{,}0 \pm 0{,}1$ \\
        $J_6$ & $ 1{,}6 \pm 0{,}1$ \\
        \bottomrule
    \end{tabular}
\end{table}

Las cinco mediciones realizadas, con las combinaciones de jinetillos,
sus masas y sus posiciones en el brazo mayor, se registran en la
Tabla~\ref{tab:datos_m1}.

\begin{table}[H]
    \centering
    \caption{Configuraciones de jinetillos registradas en el primer
             método. $n_i$: número de división del brazo mayor
             (escala de 10 divisiones).}
    \label{tab:datos_m1}
    \begin{tabular}{cccc}
        \toprule
        \textbf{Medición} & \textbf{Jinetillo} & \textbf{Masa (g)} &
        \textbf{División $n_i$} \\
        \midrule
        \multirow{2}{*}{1} & $J_1$ & $20{,}2 \pm 0{,}1$ & 1 \\
                           & $J_4$ & $ 1{,}2 \pm 0{,}1$ & 3 \\
        \midrule
        \multirow{2}{*}{2} & $J_2$ & $18{,}7 \pm 0{,}1$ & 1 \\
                           & $J_6$ & $ 1{,}6 \pm 0{,}1$ & 5 \\
        \midrule
        \multirow{2}{*}{3} & $J_4$ & $ 1{,}2 \pm 0{,}1$ & 2 \\
                           & $J_3$ & $22{,}2 \pm 0{,}1$ & 1 \\
        \midrule
        \multirow{3}{*}{4} & $J_5$ & $17{,}0 \pm 0{,}1$ & 1 \\
                           & $J_6$ & $ 1{,}6 \pm 0{,}1$ & 4 \\
                           & $J_4$ & $ 1{,}2 \pm 0{,}1$ & 2 \\
        \midrule
        \multirow{2}{*}{5} & $J_5$ & $17{,}0 \pm 0{,}1$ & 1 \\
                           & $J_6$ & $ 1{,}6 \pm 0{,}1$ & 6 \\
        \bottomrule
    \end{tabular}
\end{table}

%% ------------------------------------------------------------------
\subsection{Segundo método --- Película jabonosa con hilo}
%% ------------------------------------------------------------------

Los parámetros geométricos de la película jabonosa y el peso del tubo
inferior se presentan en la Tabla~\ref{tab:datos_m2}.

\begin{table}[H]
    \centering
    \caption{Parámetros medidos en el segundo método.}
    \label{tab:datos_m2}
    \begin{tabular}{lcc}
        \toprule
        \textbf{Magnitud} & \textbf{Símbolo} & \textbf{Valor} \\
        \midrule
        Semiancho del tubo        & $a$ & $(2{,}730 \pm 0{,}005)\ \text{cm}$ \\
        Semiancho mínimo del hilo & $b$ & $(2{,}550 \pm 0{,}005)\ \text{cm}$ \\
        Semialtura de la película & $h$ & $(2{,}780 \pm 0{,}005)\ \text{cm}$ \\
        Masa del tubo inferior    & $m$ & $(1{,}2 \pm 0{,}1)\ \text{g}$      \\
        \bottomrule
    \end{tabular}
\end{table}